% ECONOMICS_PROBLEM: {"id": "C71-12", "title": "Grand-coalition arrival in the patient limit", "jel_code": "C71", "source_id": "OP-0222"}
% !TeX program = xelatex
\documentclass[11pt,a4paper]{article}
\usepackage[margin=23mm]{geometry}
\usepackage{fontspec}
\setmainfont{Latin Modern Roman}
\usepackage{amsmath,amssymb,mathtools}
\usepackage{xcolor,enumitem,booktabs,longtable,array,xurl,hyperref}
\definecolor{muted}{HTML}{53636C}
\hypersetup{colorlinks=true,urlcolor=blue,linkcolor=blue}
\setlength{\parindent}{0pt}
\setlength{\parskip}{5pt}
\newcommand{\R}{\mathbb R}
\newcommand{\E}{\mathbb E}
\newcommand{\N}{\mathbb N}
\newcommand{\ind}{\mathbf 1}
\newcommand{\Var}{\operatorname{Var}}
\newcommand{\Cov}{\operatorname{Cov}}
\newcommand{\supp}{\operatorname{supp}}
\DeclareMathOperator*{\argmin}{arg\,min}
\DeclareMathOperator*{\argmax}{arg\,max}
\newcommand{\entrymeta}[3]{{\sffamily\small\color{muted}\textbf{\texttt{#1}}\quad|\quad #2\par #3\par}}
\newcommand{\Problemref}[1]{\texttt{#1}}
\newcommand{\JELref}[2]{\texttt{#2} (JEL \texttt{#1})}
\begin{document}
\section*{Grand-coalition arrival in the patient limit}
\textbf{Problem ID:} \texttt{C71-12}\quad \textbf{JEL:} \texttt{C71}\par
% BEGIN ECONOMICS STATEMENT
\label{problem:OP-0222}
% GAME_THEORY_PROBLEM: {"local_id": "GT-092", "original_number": 92, "kind": "P", "entry_kind": "statement_only", "published": false, "review_required": true, "category": "game_theory"}
\label{game:GT-092}
{\sffamily\small\color{muted}
\noindent\textbf{ID:} GT-092. \textbf{Category:} Dynamic coalition formation.
\textbf{Status:} P for the patient-limit question; the catalog's counterexample claim requires a version correction.
\par}
\subsection*{Definitions and mathematical statement}
Fix a finite player set $N$, $|N|\geq2$. Let $\mathcal X$ be the set of laminar families of subsets of $N$ of size at least two: any two members are disjoint or one contains the other. For $x\in\mathcal X$, let $P(x)$ comprise its maximal members and the singleton players outside their union. Put
$\mathcal G=\{x:N\in x\}$. Payoffs are $\pi_i:\mathcal X\to\mathbb R$ and $V(x)=\sum_i\pi_i(x)$.
Require the source's Assumptions 2.8 and 2.9:
\[
 V(h)=V^*\ (h\in\mathcal G),\quad
 \gamma=\min_{x\notin\mathcal G}(V^*-V(x))>0,
\]
and, for each $K\in x\cap P(x)$, positive shares $s_i(x,K)$ summing to one over $i\in K$, with
\[
 \Delta(x,K)=\sum_{j\in K}\bigl(\pi_j(x)-\pi_j(x\setminus\{K\})\bigr),\qquad
 \pi_i(x)=\pi_i(x\setminus\{K\})+s_i(x,K)\Delta(x,K).
\]

A move has a target $\tau(m)$ and a nonempty relevant player set $R(m)$. A merge chooses at least two blocks $\mathcal A\subseteq P(x)$, sets $K=\bigcup\mathcal A$, and has target $x\cup\{K\}$ and $R(m)=K$. A unilateral termination chooses $K\in x$ and $i\in K$, and has target
$x\setminus\{L\in x:K\subseteq L\}$ and $R(m)=\{i\}$. A stay by $i$ has target $x$ and $R(m)=\{i\}$. Write $M(x)$ for moving moves and $M^+(x)$ for all these moves.

For $\delta\in(0,1)$ and a stochastic transition matrix $p$ on $\mathcal X$, define
\[
 \ell_i=(1-\delta)(I-\delta p)^{-1}\pi_i,\qquad
 \ell_i^+(x)=\sum_y p_{xy}\ell_i(y).
\]
At $x$, rank targets by $\ell_i$, breaking equal values in favour of $x$; denote this total preorder by $\succeq_i^x$ and its strict part by $\succ_i^x$. A move $m$ is profitable if
$\ell_i(\tau(m))\geq\ell_i^+(x)$ for every $i\in R(m)$. It is dominated if some $m'\in M^+(x)$ satisfies $R(m')\subseteq R(m)$ and
$\tau(m')\succ_i^x\tau(m)$ for every $i\in R(m')$.
Let $C(x)$ be the profitable undominated moves, and
\[
 F(x)=\{m\in C(x):\exists i\in N\ \forall m'\in C(x),\
                  \tau(m)\succeq_i^x\tau(m')\}.
\]
The matrix $p$ is an equilibrium if, at every $x$,
\begin{align*}
 p_{xy}>0, y\neq x&\ \Longrightarrow\ \exists m\in C(x):\tau(m)=y,\\
 C(x)\cap M(x)\neq\varnothing&\ \Longrightarrow\ p_{xx}<1,\\
 m\in F(x)&\ \Longrightarrow\ p_{x,\tau(m)}>0.
\end{align*}
For the resulting Markov chain let $T_{\mathcal G}=\inf\{t\geq0:X_t\in\mathcal G\}$.

Does there exist \emph{one fixed} payoff model satisfying the assumptions, a sequence $\delta_k\uparrow1$, equilibrium matrices $p_k$, and states $x_k$ such that
\[
 \mathbb P_{x_k}^{p_k}(T_{\mathcal G}=\infty)>0\quad\text{for every }k?
\]
Equivalently, is it true that every fixed model has a threshold $\bar\delta<1$ beyond which every equilibrium arrives at $\mathcal G$ almost surely from every initial state?

\subsection*{Short English statement}
Can a fixed coalition-formation model sustain nonarrival for equilibria with discount factors arbitrarily close to one?

\subsection*{Variants and scope boundaries}
The fixed-model quantifier is stronger than allowing payoffs or the number of players to change with $k$. A model-dependent patience threshold is weaker than a universal threshold. The question uses unilateral termination and sharing of per-period surplus. Sharing discounted surplus, unanimous termination, or deleting only a nested agreement while preserving its ancestors changes the problem.

\subsection*{Source relationship and status}
The definitions follow the inspected v2's Sections 2.1--2.4, and the patient-limit target follows its Section 8. Remark 6.7 and the heading of Proposition 6.8 explicitly state that the earlier deterministic-cycle construction is not an equilibrium under revised Definition 2.14, under either termination rule. Consequently the catalog's claim of established failure for $\delta\leq3/4$ cannot be transferred to these definitions. The formulation above fixes the primitives, making one precise patient-limit question; the source's broader phrase does not itself distinguish all uniformity variants.

\subsection*{References}
\begingroup\footnotesize\raggedright
\noindent [S1] Jobst Heitzig, \emph{Does the grand coalition form? Persistence, arrival, and the role of the sharing rule in a dynamic process of nested binding agreements} (2026), arXiv:2608.17766v2. Inspected locators: Definitions 2.1--2.4, Assumptions 2.8--2.9, Definitions 2.14--2.15, Remark 6.7 and Proposition 6.8, Section 8, Open problems (2).
\url{https://arxiv.org/html/2608.17766v2}
\par\endgroup

\subsection*{Common conventions: Source mathematical conventions}

\textbf{Finite games.} For a finite set $A$, $\Delta(A)$ is its probability
simplex. Players choose independent mixed strategies in Nash equilibrium;
correlation is allowed only when explicitly part of the solution concept.
In a game with payoff functions $u_i$ and strategy profile $x$, an additive
$\varepsilon$ Nash equilibrium satisfies
\[
 u_i(x)\geq u_i(x_i',x_{-i})-\varepsilon
 \quad\text{for every player $i$ and every admissible deviation $x_i'$}.
\]
For numerical approximation, payoffs are normalized as locally specified.
An $\varepsilon$ payoff residual is distinct from distance at most
$\varepsilon$ to the set of exact equilibria. Point convergence, convergence
of empirical play, and convergence in distance to an equilibrium set are
also distinct.

\textbf{Computational representations.} Unless stated otherwise, a complexity
question takes rational data with binary encoding; polynomial time is measured
in total bit length. A fixed-accuracy polynomial algorithm is not necessarily
polynomial in $1/\varepsilon$, and neither is necessarily polynomial in
$\log(1/\varepsilon)$. Succinct games and value, demand, or sampling oracles
must declare their interfaces. Exact real solutions require an explicit
representation; an unspecified list of arbitrary real numbers is not a
finite computational output. A claimed algorithmic barrier is meaningful
only for its specified model and reductions.

\textbf{Dynamic games.} Histories, information, observation, and allowable
randomization are part of the model. A stationary strategy depends only on
the current state; a positional strategy in a deterministic graph game
chooses one action at each controlled vertex. Nash equilibrium at an initial
state and equilibrium after every history are different requirements.
An expectation of a pathwise $\liminf$ payoff, a $\liminf$ of expected averages,
a limiting expected average, and a uniform finite-horizon equilibrium are
different criteria. None is silently substituted for another.

\textbf{Allocation and mechanisms.} A complete allocation assigns every item
exactly once unless otherwise stated. Zero-valued goods, negative values,
capacity constraints, feasibility, lotteries, and copies of goods affect
fairness definitions and are specified locally. Dominant-strategy incentive
compatibility, Bayesian incentive compatibility, universal truthfulness,
and truthfulness in expectation impose different inequalities. Whenever
an objective ratio has zero denominator, its convention is stated or those
instances are excluded.

\textbf{Infinite-dimensional models.} Expectations, integrals, stochastic
equations, and optimization objectives require their local regularity,
integrability, filtrations, and admissible domains. A backward--forward
mean-field system is a boundary-value problem; it is not an initial-value
semiflow without an additional construction. Long-time, large-population,
and vanishing-noise limits require an order of limits and a convergence
topology. If a minimum need not be attained, the objective is its infimum
or supremum and the approximation requirement is stated separately.
% END ECONOMICS STATEMENT
\end{document}
